\documentclass[10pt,a4paper]{article}
\usepackage{amsmath,amssymb,amsthm,graphicx}
\usepackage{amsmath,amssymb,amsthm,hyperref,xfrac}
\providecommand{\ie}{i.e.\ }
\providecommand{\skpt}{\smallskip}
\usepackage[scr=rsfs,cal=euler]{mathalfa}
\multlinegap0pt

\hyphenation{begin decom-po-si-tion sequence}
\newenvironment{enumeratea}
{\bgroup\def\theenumi{\alph{enumi}}\begin{enumerate}}
{\end{enumerate}\egroup}

\newcommand{\ba}{{\boldsymbol a}}
\newcommand{\bb}{{\boldsymbol b}}
\newcommand{\bd}{{\boldsymbol d}}
\newcommand{\bu}{{\boldsymbol u}}
\newcommand{\bx}{{\boldsymbol x}}
\newcommand{\bmy}{{\boldsymbol y}}

\newcommand{\bD}{{\boldsymbol D}}
\newcommand{\bL}{{\boldsymbol L}}
\newcommand{\bP}{{\boldsymbol P}}
\newcommand{\bV}{{\boldsymbol V}}
\newcommand{\bX}{{\boldsymbol X}}
\newcommand{\bZ}{{\boldsymbol Z}}

\newcommand{\Psfrac}[2]{(\sfrac{#1}{#2})}
\newcommand{\psfrac}[2]{\sfrac{(#1)}{#2}}
\newcommand{\bigO}[1]{O(#1)}

\usepackage{mathtools}
\usepackage{bbm}

\newcommand{\NN}{\mathbb{N}}
\newcommand{\ZZ}{\mathbb{Z}}
\newcommand{\RR}{\mathbb{R}}
\newcommand{\CC}{\mathbb{C}}
\newcommand{\PP}{\mathbb{P}}
\newcommand{\EE}{\mathbb{E}}

\newcommand{\boldalpha}{\boldsymbol{\alpha}}
\newcommand{\bolddelta}{\boldsymbol{\delta}}
\newcommand{\boldrho}{\boldsymbol{\rho}}
\newcommand{\dee}{\,\mathrm{d}}
\newcommand{\one}{\mathbbm{1}}
\newcommand{\BC}{\mathcal{B}}
\newcommand{\EC}{\mathcal{E}}
\newcommand{\RS}{\mathscr{R}}
\newcommand{\prob}[1]{\PP[ #1 ]}
\newcommand{\bigoo}[1]{O\!\left(#1 \right)}
\newcommand{\abs}[1]{\left\lvert #1 \right\rvert}
\newcommand{\floor}[1]{\left\lfloor #1 \right\rfloor}
\newcommand{\riemint}[4]{\int_{#1}^{#2}\! #3 \, \mathrm{d}#4}
\newcommand{\maxnorm}[1]{\left\| #1 \right\|_{\infty}}

\renewcommand{\Re}{\textup{Re}}
\renewcommand{\Im}{\textup{Im}}

\newcommand{\dtv}{\mathrm{d}_{\mathrm{TV}}}
\newcommand{\dw}{\mathrm{d}_{\mathrm W}}

\let\epsilon\varepsilon

\newcommand{\pextra}{P_{\mathtt{extra}}}
\newcommand{\parratia}{P_{\mathtt{Arratia}}}

\newcommand{\simplex}{\Delta^{k-1}}

\newcommand{\bg}{\big}
\newcommand{\bgg}{\Big}
\newcommand{\bggg}{\bigg}

\newcommand{\Rnt}{R_{\mathrm{NT}}}
\newcommand{\Rprob}{R_{\mathrm{PD}}}
\newcommand{\ECnt}{\EC_{\mathrm{NT}}}
\newcommand{\ECprob}{\EC_{\mathrm{PD}}}

\newcommand{\Bnt}{\mathcal{B}_{\mathrm{NT}}}
\newcommand{\Bprob}{\mathcal{B}_{\mathrm{PD}}}

\DeclareMathOperator{\vard}{Var}
\DeclareMathOperator{\Dir}{Dir}

\theoremstyle{plain}
\newtheorem{cor}{Corollary}[section]
\newtheorem{lem}[cor]{Lemma}
\newtheorem{prop}[cor]{Proposition}
\newtheorem{conj}[cor]{Conjecture}
\newtheorem{theo}{Theorem}

\theoremstyle{definition}
\newtheorem{example}{Example}
\newtheorem{deff}[cor]{Definition}
\newtheorem{rem}{Remark}[section]
\newtheorem*{rems}{Remark}
\newtheorem*{remss}{Remarks}

\makeatletter
\ams@newcommand{\multiint}[1]{\DOTSI\protect\MultiIntegral{#1}}
\renewcommand{\MultiIntegral}[1]{%
\edef\ints@c{\noexpand\intop
\ifnum#1=\z@\noexpand\intdots@\else\noexpand\intkern@\fi
\replicate{#1-2}{\noexpand\intop\noexpand\intkern@}%
\noexpand\intop
\noexpand\ilimits@
}%
\futurelet\@let@token\ints@a
}
\makeatother
\ExplSyntaxOn
\cs_new:Npn \replicate #1 #2 { \prg_replicate:nn { #1 } { #2 } }
\ExplSyntaxOff
\title{Reduction of the boundary-event estimate to a two-parameter Poisson integral bound}
\author{Tony Haddad, Dimitris Koukoulopoulos}
\date{doi:10.5802/jep.317; CC BY 4.0}
\begin{document}
\maketitle
\noindent\textit{Editorial note.} The counted statement is the author's labelled lemma `lem:Bprob reduction 2` (source lines 1563-1569); the proof environment begins on the next line after \textbackslash{}end\{lem\} apart from one blank line, so the pairing is by adjacency.\par\medskip
\noindent\textit{Editorial note.} Three context spans are retained: the author's own Dirichlet law proposition whose proof supplies the Gamma representation, the uniform Beta integral estimate invoked at the last step (also counted separately as its own item), and the appendix passage containing Mecke's equation.\par\medskip
\noindent\textit{Editorial note.} No counter adjustment is retained: the publisher numbering is not reproduced and the item prints with the wrapper's own sequential numbering.\par\medskip
\noindent\textit{Editorial note.} Every label used inside the retained text is defined inside one of the retained spans, so the package declares no external reference.\par\medskip
\noindent\textit{Editorial scope.} One original result of the pinned publisher article, counted once: the two-parameter reduction bound for the boundary event (author label lem:Bprob reduction 2).\par\medskip
\noindent\textit{Allowed standard prerequisites (editorial).}\begin{itemize}
\item Poisson process representation of the Dirichlet components. Splitting a Poisson process of intensity e\textasciicircum{}\{-x\}/x dx into two independent parts with intensities alpha e\textasciicircum{}\{-x\}/x dx and (1-alpha) e\textasciicircum{}\{-x\}/x dx gives S\_A \textasciitilde{} Gamma(alpha,1) and S\_B \textasciitilde{} Gamma(1-alpha,1), with S\_A/S Dirichlet distributed.. It is the representation introduced at source lines 1572-1573 that replaces the expectation E by an expectation over the two Poisson processes.
\item Mecke's equation. For a Poisson process, the expected sum over ordered pairs of a functional of the points equals the corresponding integral against the intensity measure.. It is the step at source lines 1595-1601 that turns the double sum in E\_A into the displayed four-fold integral.
\item Beta integral estimate. sin(pi alpha) times the integral of t\textasciicircum{}\{alpha-1\}(1-t)\textasciicircum{}\{-alpha\} over an interval of length about (10 lambda + mu)/log x is bounded by (lambda+mu) times Delta.. It is applied at source line 1627 to close the bound for E\_A, and the same bound covers E\_B.
\end{itemize}\medskip
\begin{prop}[Donnelly--Tavaré \cite{DonnellyTavare87}]
\label{prop:probabilistic dirichlet law}
Let $\bV = (V_1, V_2, \ldots)$ be a Poisson--Dirichlet process, let $\boldalpha = (\alpha_1, \ldots, \alpha_k) \in \simplex$ such that $\alpha_i > 0$ for all $1 \le i \le k$, and let $(C_i)_{i \ge 1}$ be a sequence of i.i.d.~random variables (also independent of $\bV$) with $\PP[C_i = j] = \alpha_j$ for all $i \ge 1$. Then the random vector
\[
\Bigl(\sum_{i:\ C_i = 1} V_i, \ldots, \sum_{i:\ C_i = k} V_i\Bigr)
\]
is distributed according to $\Dir(\boldalpha)$.
\end{prop}

\begin{lem}\label{lem:beta estimate}
Uniformly for $u\in[0,1]$, $\alpha \in (0, 1)$ and $\delta > 0$, we~have
\[
\sin(\pi\alpha) \int_{[u-\delta,u+\delta]\cap[0,1]} \frac{\dee t}{t^{1-\alpha}(1-t)^\alpha} \ll \frac{\delta}{(u+\delta)^{1-\alpha}(1-u+\delta)^\alpha}.
\]
In particular, if $\delta\ge1/\log x$, then
\[
\sin(\pi\alpha) \int_{[u-\delta,u+\delta]\cap[0,1]} \frac{\dee t}{t^{1-\alpha}(1-t)^\alpha} \ll \frac{\delta\log x}{(1+u\log x)^{1-\alpha}(1+(1-u)\log x)^\alpha}.
\]
\end{lem}


\subsection{Poisson Processes}

The following definition and propositions are borrowed from Kingman's book on Poisson processes \cite{Kingman93}. Let $(S, \mathcal S)$ be a measurable space with~$S$ being a subset of $\RR^d$ for some $d \ge 1$. A \textit{Poisson process} on a state space $S$ with mean measure $\mu$ is a random countable subset $\Pi \subseteq S$ such that:
\begin{itemize}
\item for any disjoint measurable subsets $A_1, \ldots, A_n$ of $S$, the random variables \hbox{$\#(\Pi \cap A_1)$, \dots, $\#(\Pi \cap A_n)$} are independent;

\item the random variable $\#(\Pi \cap A)$ is a Poisson random variable of parameter $\mu(A)$ for any $A\subset S$ measurable.
\end{itemize}
We can see $\Pi$ as an element of the measurable space $(\Omega_S, \mathcal F_S)$ where $\Omega_S$ is the set of countable subsets of $S$ and $\mathcal F_S$ is the smallest $\sigma$-algebra for which the map $\Pi \mapsto \#\{\Pi \cap\nobreak B\}$ is measurable for all $B \in \mathcal S$. If~$\mu$ has no atoms, meaning no singleton with positive probability, and is $\sigma$-finite, meaning that $S$ is a countable union of measurable sets with finite measure, then a Poisson process with mean measure $\mu$ always exists (see \cite{Kingman93}, the existence theorem in section 2.5 for a proof). If~$\mu$ is absolutely continuous with respect to the Lebesgue measure, then the function \hbox{$\lambda \colon S \to \RR_{\ge 0}$} such that $\mu(A) = \riemint{A}{}{\lambda(x)}{x}$ for all measurable subsets $A \subseteq S$ is called the \textit{intensity} of the Poisson process. Here are a few important propositions about Poisson processes that we will use in the paper and proved in \cite{Kingman93}:

\begin{prop}[Mapping theorem, {\cite[\S2.3]{Kingman93}}]
\label{prel:map}
If $\Pi$ is a Poisson process with mean measure $\mu$ on $S$, and $f:S \to T$ is a measurable function such that $\mu^*(B) \coloneqq \mu(f^{-1}(B))$ has no atoms, then $f(\Pi)$ is a Poisson process on $T$ with measure $\mu^*$.
\end{prop}

\begin{prop}[Coloring theorem, {\cite[\S5.1]{Kingman93}}]
\label{prel:color}
If $\Pi$ is a Poisson process with mean measure $\mu$ on $S$, and the points are randomly colored with $k$ colors such that the probability of a point receiving the color $i$ is $p_i$, and such that the color of a point is independent of different points and of the position of the point. Let $\Pi_i$ be the subset of~$\Pi$ with color $i$. Then all the $\Pi_i$ are independent Poisson processes with mean measures $\mu_i = p_i\mu$.
\end{prop}

\begin{prop}[Campbell's theorem, {\cite[\S3.2]{Kingman93}}]
\label{prel:campbell}
Let $\Pi$ be any Poisson process on $S$ with mean measure $\mu$. Let $f\colon S \to \RR$ be a measurable function. Then
\[
\Sigma = \sum_{X \in \Pi} f(X)
\]
is absolutely convergent almost surely if, and only if,
\[
\riemint{S}{}{\min\{f, 1\}}{\mu} < \infty.
\]
If this condition holds, then
\[
\EE[e^{s\Sigma}] = \exp\left(\riemint{S}{}{(e^{sf} - 1)}{\mu}\right)
\]
for any complex $s$ for which the integral converges. Moreover,
\[
\EE[\Sigma] = \riemint{S}{}{f}{\mu}
\]
if the integral converges. In~the case where it converges, we~also have
\[
\vard[\Sigma] = \riemint{S}{}{f^2}{\mu}.
\]
\end{prop}

Many probabilities or expectations that involve Poisson processes in this paper can be reformulated as
\[
\EE \sum_{X \in \Pi} f(\Pi\setminus \{X\}, X).
\]
In these cases, there is a generalization of the formula for $\EE[\Sigma]$ in Campbell's theorem, called the Mecke equation, allowing us to compute these objects:

\begin{prop}[Mecke equation, {\cite[Th.\,4.5]{LastPenrose18}}]
\label{prel:mecke}
Let $\Pi$ be a Poisson process on $S$ with a $\sigma$-finite mean measure $\mu$, and let $f: \Omega_S \times S^k \to [0, \infty)$ be measurable. Then we have
\begin{multline*}
\EE \sum_{\substack{X_1, \ldots, X_k \\ \text{all distinct}}} f(\Pi\setminus \{X_1, \ldots, X_k\}; X_1, \ldots, X_k) \\[-10pt]
= \int_S \cdots \int_S \EE[f(\Pi; x_1, \ldots, x_k)] \, \mathrm d\mu(x_1)\, \cdots\, \mathrm d\mu(x_k).
\end{multline*}
\end{prop}

\begin{lem}
\label{lem:Bprob reduction 2} Uniformly for $\mu\ge1$ and $\lambda\in(0,0.1\log x]$, we~have
\[
\EE\bgg[ \one_{| Z_i - u |\le \sfrac{\mu}{\log x} } \cdot \one_{G(\lambda) \ge4} \cdot G(\lambda)^2 \bgg]
\ll (\lambda+\mu)\cdot \Delta.
\]
\end{lem}

\begin{proof} Let us call $E$ the quantity we seek to bound from above.
Consider two independent Poisson processes $(A_j)_{j \ge 1}$ with intensity $\alpha\Psfrac{e^{-x}}{x} \dee x$ and $(B_j)_{j \ge 1}$ with intensity $(1-\alpha)\Psfrac{e^{-x}}{x} \dee x$. Hence, the union of the two processes is a new Poisson process of intensity $\Psfrac{e^{-x}}{x}\dee x$.
The sum $S_A \coloneqq \sum_{j\ge1} A_j$ has distribution $\text{Gamma}(\alpha, 1)$ and the sum $S_B \coloneqq \sum_{j\ge1} B_j$ has distribution $\text{Gamma}(1-\alpha, 1)$. If~we also set $S\coloneqq S_A+S_B$, then, as we saw in the proof of Proposition \ref{prop:probabilistic dirichlet law}, we~have
\[
E
=\EE\bgg[ \one_{| S_A/S- u |\le \sfrac{\mu}{\log x} } \cdot \one_{G_A(\lambda)+G_B(\lambda)\ge 4}\cdot \bg( G_A(\lambda)+G_B(\lambda) \bg)^2 \bgg],
\]
where $G_A(\lambda) \!=\! \sum_{j\ge1} \one_{A_j\log x\in(\lambda S,2\lambda S]}$ and $G_B(\lambda)\!=\! \sum_{j\ge1} \one_{B_j\log x\in(\lambda S,2\lambda S]}$. We~know $G_A(\lambda)+G_B(\lambda)\ge4$. Hence, if $G_A(\lambda)\ge G_B(\lambda)$, then we must have \hbox{$G_A(\lambda)\ge 2$} and $G_A(\lambda)\le G_A(\lambda)^2/2$, whence $G_A(\lambda) \le 2\mathop{\sum\sum}_{k>j\ge1} \one_{A_k\log x,\, A_j\log x\in (\lambda S,2\lambda S]}$. So, we~find
\[
\begin{split}
\bg( G_A(\lambda)+G_B(\lambda) \bg)^2 \le 4G_A(\lambda)^2
&= 4 G_A(\lambda) + 8 \mathop{\sum\sum}_{k>j\ge1} \one_{A_k\log x,\, A_j\log x\in (\lambda S,2\lambda S]} \\
&\le 16 \mathop{\sum\sum}_{k>j\ge1} \one_{A_k\log x,\, A_j\log x\in (\lambda S,2\lambda S]}.
\end{split}
\]
The analogous inequality also holds when $G_B(\lambda)\ge G_A(\lambda)$, with the roles of $A$ and~$B$ reversed. We~conclude that
\[
E
\le 16 (E_A+E_B),
\]
where
\[
E_A \coloneqq \EE\bgg[\textstyle \one_{|S_A/S - u |\le \sfrac{\mu}{\log x} } \mathop{\sum\sum}_{k>j\ge1} \one_{A_k\log x,\, A_j\log x\in (\lambda S,2\lambda S]} \bgg]
\]
and $E_B$ is defined analogously. Using Mecke's equation (cf. Proposition \ref{prel:mecke}), we~find
\[
\begin{split}
E_A &= \alpha^2\iint\limits_{2a_1>a_2>a_1>0}
\EE\bgg[\textstyle \one_{| \frac{a_1+a_2+S_A}{a_1+a_2+S} - u |\le \frac{\mu}{\log x} } \prod_{j=1}^2 \one_{a_j\log x \in (\lambda (a_1+a_2+S),2\lambda (a_1+a_2+S)]} \bgg] \\[-10pt]
&\hspace*{9.5cm} \times \frac{e^{-a_1-a_2}}{a_1a_2}\dee a_1\dee a_2 \\[6pt]
&= \frac{\alpha^2}{\Gamma(\alpha)\Gamma(1-\alpha)} \mkern-50mu \iiiint\limits_{
\substack{
2a_1>a_2>a_1>0,\ s_1,s_2>0 \\
| \frac{a_1+a_2+s_1}{a_1+a_2+s_1+s_2} - u |\le \frac{\mu}{\log x} \\
a_j\log x \in (\lambda (a_1+a_2+s_1+s_2),2\lambda (a_1+a_2+s_1+s_2)] \ (j=1,2)
} }
\mkern-50mu\frac{e^{-a_1-a_2-s_1-s_2}}
{a_1a_2s_1^{1-\alpha}s_2^\alpha}\dee a_1\dee a_2 \dee s_1\dee s_2,
\end{split}
\]
where we used that the $a_j$'s lie in the same dyadic interval to deduce that $a_2<2a_1$. We~make the change of variables $t=s_1/(s_1+s_2)$ and $s=s_1+s_2$. Since $\lambda/\log x\le 0.1$ and $a_1<a_2<2a_1$, the conditions $a_j\log x \in (\lambda (a_1+a_2+s),2\lambda (a_1+a_2+s)]$ for $j=1,2$ imply that $a_j\log x \in(\lambda s,5\lambda s]$. Knowing also that $ |\frac{a_1+a_2+s_1}{a_1+a_2+s} - u |\le \frac{\mu}{\log x}$, we~find $|t-u|=|\frac{s_1}{s_1+s_2} - u|\le \frac{10\lambda+\mu}{\log x}$. Finally, we~use Euler's reflection formula to write $\Gamma(\alpha)\Gamma(1-\alpha)=\pi/\sin(\pi\alpha)$. We~conclude that
\[
E_A \le \frac{\sin(\pi \alpha)}{\pi}
\iiiint\limits_{
\substack{
a_2>a_1>0,\ s>0,\ t\in[0,1] \\
| t-u|\le (10\lambda+\mu)/\log x\\
a_j\log x \in (\lambda s, 5\lambda s] \ (j=1,2)
} }
\frac{e^{-a_1-a_2-s}}
{a_1a_2t^{1-\alpha}(1-t)^\alpha}\dee a_1\dee a_2 \dee s\dee t.
\]
For every fixed value of $s$, the integral over $a_1$ and $a_2$ is $\le(\log 5)^2$, since $\int_w^{5w}\frac{e^{-a}}{a}\dee a \le \log5$ for any $w>0$. We~also have $\int_0^\infty e^{-s}\dee s=1$. We~thus conclude that
\[
E_A\le \frac{(\log 5)^2\sin(\pi \alpha) }{\pi} \int\limits_{\substack{t\in[0,1] \\ |t-u|\le (10\lambda+\mu)/\log x}} \frac{\dee t}{t^{1-\alpha}(1-t)^\alpha}.
\]
Using Lemma \ref{lem:beta estimate} shows that $E_A\ll (\lambda+\mu)\Delta$. The same estimate holds for $E_B$ too, thus completing the proof of the lemma.
\end{proof}

\begin{thebibliography}{99}
\bibitem{Arratia98} R. Arratia - “On the central role of scale invariant Poisson processes on (0,∞)”, in Microsurveys in discrete probability (Princeton, NJ, 1997), DIMACS Ser. Discrete Math. Theoret. Comput. Sci., vol. 41, American Mathematical Society, Providence, RI, 1998, p. 21-41
\bibitem{Arratia02} R. Arratia - “On the amount of dependence in the prime factorization of a uniform random integer”, in Contemporary combinatorics, Bolyai Soc. Math. Stud., vol. 10, János Bolyai Math. Soc., Budapest, 2002, p. 29-91
\bibitem{ABT06} R. Arratia, A. D. Barbour \& S. Tavaré - “A tale of three couplings: Poisson-Dirichlet and GEM approximations for random permutations”, Combin. Probab. Comput. 15 (2006) no. 1-2, p. 31-62
\bibitem{BareikisManstavicius07} G. Bareikis \& E. Manstavičius - “On the DDT theorem”, Acta Arith. 126 (2007) no. 2, p. 155-168
\bibitem{Bill72} P. Billingsley - “On the distribution of large prime divisors”, Period. Math. Hungar. 2 (1972), p. 283-289
\bibitem{delaBretecheTenenbaum16} R. de la Bretèche \& G. Tenenbaum - “Sur les processus arithmétiques liés aux diviseurs”, Adv. in Appl. Probab. 48 (2016) no. A, p. 63-76
\bibitem{DDT79} J.-M. Deshouillers, F. Dress \& G. Tenenbaum - “Lois de répartition des diviseurs. I”, Acta Arith. 34 (1979) no. 4, p. 273-285 (loose errata)
\bibitem{DonnellyTavare87} P. Donnelly \& S. Tavaré - “The population genealogy of the infinitely-many neutral alleles model”, J. Math. Biol. 25 (1987) no. 4, p. 381-391
\bibitem{Feng10} S. Feng - The Poisson-Dirichlet distribution and related topics. Models and asymptotic behaviors, Probability and its Appl., Springer, Heidelberg, 2010
\bibitem{JKB} N. L. Johnson, S. Kotz \& N. Balakrishnan - Discrete multivariate distributions, Wiley Series in Probab. and Stat.: Applied Probab. and Stat., John Wiley \& Sons, Inc., New York, 1997
\bibitem{Kingman93} J. F. C. Kingman - Poisson processes, Oxford Studies in Probab., vol. 3, The Clarendon Press, Oxford University Press, New York, 1993
\bibitem{KotzBal00} S. Kotz, N. Balakrishnan \& N. L. Johnson - Continuous multivariate distributions. Vol. 1. Models and applications, Wiley Series in Probab. and Stat.: Applied Probab. and Stat., Wiley-Interscience, New York, 2000
\bibitem{Koukoulo19} D. Koukoulopoulos - The distribution of prime numbers, Graduate Studies in Math., vol. 203, American Mathematical Society, Providence, RI, 2019
\bibitem{LastPenrose18} G. Last \& M. Penrose - Lectures on the Poisson process, Institute of Math. Stat. Textbooks, vol. 7, Cambridge University Press, Cambridge, 2018
\bibitem{Ken22} S.-K. Leung - “Dirichlet law for factorisation of integers, polynomials and permutations”, Math. Proc. Cambridge Philos. Soc. 175 (2023) no. 3, p. 649-676
\bibitem{KenWeb24} S.-K. Leung - “Dirichlet law with arbitrary parameters for factorizations of integers” (2025), Retrieved on 24 August 2025 from https://sites.google.com/view/sunkaileung/miscellaneous
\bibitem{MV07} H. L. Montgomery \& R. C. Vaughan - Multiplicative number theory. I. Classical theory, Cambridge Studies in Advanced Math., vol. 97, Cambridge University Press, Cambridge, 2007
\bibitem{NyandwiSmati13} S. Nyandwi \& A. Smati - “Distribution laws of pairs of divisors”, Integers 13 (2013), article ID A13, 13 pages
\bibitem{Tenenbaum00} G. Tenenbaum - “A rate estimate in Billingsley’s theorem for the size distribution of large prime factors”, Q. J. Math. 51 (2000) no. 3, p. 385-403
\end{thebibliography}
\end{document}
